\documentclass[a4paper,14pt,oneside,final]{report}
% \usepackage{anyfontsize}
\usepackage[utf8]{inputenc}
\usepackage[russian]{babel}
\usepackage[T2A]{fontenc}
% \usepackage{anyfontsize}
\usepackage{graphicx}
\usepackage{float}
\DeclareGraphicsExtensions{.pdf,.png,.jpg,.jpeg}
\graphicspath{{png/}}
\usepackage{tocloft}
\renewcommand{\cfttoctitlefont}{\normalfont\Large\bfseries\centering}
\renewcommand{\contentsname}{\MakeUppercase{Содержание}}
\setlength{\cftbeforetoctitleskip}{-1em}
\setlength{\cftaftertoctitleskip}{1em}
%\usepackage{pscyr}
%\renewcommand{\rmdefault}{cmr}
\usepackage[final,hidelinks]{hyperref}
\usepackage[numbers,square,sort&compress]{natbib}
%\usepackage{times}
%\usefont{T2A}{ftm}{m}{sl}
\setlength{\bibsep}{0em}
\PassOptionsToPackage{hyphens}{url}
\usepackage{pgf}
\usepackage{enumitem} % подключи в преамбуле
\usepackage{tikz}
\usetikzlibrary{positioning, fit, shapes.geometric, shapes.symbols, shadows, arrows.meta}
\usepackage{float}
\usepackage{listings} % Для листингов кода
\usepackage{xcolor}
\usepackage{hyphenat}
\lstset{
	language=Python, % Указываем язык Python
	basicstyle=\small\ttfamily, % Шрифт для кода
	keywordstyle=\color{blue}\bfseries, % Стиль ключевых слов
	stringstyle=\color{red}, % Стиль строк
	commentstyle=\color{green!50!black}\itshape, % Стиль комментариев
	numbers=left, % Номера строк слева
	numberstyle=\tiny, % Размер шрифта для номеров строк
	stepnumber=1, % Шаг нумерации строк
	numbersep=5pt, % Отступ номеров строк
	showspaces=false, % Не показывать пробелы
	showstringspaces=false, % Не показывать пробелы в строках
	frame=single, % Рамка вокруг кода
	breaklines=true, % Перенос длинных строк
	breakatwhitespace=true, % Перенос только на пробелах
	tabsize=4, % Размер табуляции
	literate={á}{{\'a}}1 {é}{{\'e}}1 {í}{{\'i}}1 {ó}{{\'o}}1 {ú}{{\'u}}1
	{А}{{\cyrchar\CYRA}}1 {Б}{{\cyrchar\CYRB}}1 {В}{{\cyrchar\CYRV}}1
	{Г}{{\cyrchar\CYRG}}1 {Д}{{\cyrchar\CYRD}}1 {Е}{{\cyrchar\CYRE}}1
	{Ё}{{\cyrchar\CYRYO}}1 {Ж}{{\cyrchar\CYRZH}}1 {З}{{\cyrchar\CYRZ}}1
	{И}{{\cyrchar\CYRI}}1 {Й}{{\cyrchar\CYRISHRT}}1 {К}{{\cyrchar\CYRK}}1
	{Л}{{\cyrchar\CYRL}}1 {М}{{\cyrchar\CYRM}}1 {Н}{{\cyrchar\CYRN}}1
	{О}{{\cyrchar\CYRO}}1 {П}{{\cyrchar\CYRP}}1 {Р}{{\cyrchar\CYRR}}1
	{С}{{\cyrchar\CYRS}}1 {Т}{{\cyrchar\CYRT}}1 {У}{{\cyrchar\CYRU}}1
	{Ф}{{\cyrchar\CYRF}}1 {Х}{{\cyrchar\CYRH}}1 {Ц}{{\cyrchar\CYRC}}1
	{Ч}{{\cyrchar\CYRCH}}1 {Ш}{{\cyrchar\CYRSH}}1 {Щ}{{\cyrchar\CYRSHCH}}1
	{Ъ}{{\cyrchar\CYRHRDSN}}1 {Ы}{{\cyrchar\CYRY}}1 {Ь}{{\cyrchar\CYRSFTSN}}1
	{Э}{{\cyrchar\CYRE}}1 {Ю}{{\cyrchar\CYRYU}}1 {Я}{{\cyrchar\CYRYA}}1
	{а}{{\cyrchar\cyra}}1 {б}{{\cyrchar\cyrb}}1 {в}{{\cyrchar\cyrv}}1
	{г}{{\cyrchar\cyrg}}1 {д}{{\cyrchar\cyrd}}1 {е}{{\cyrchar\cyre}}1
	{ё}{{\cyrchar\cyryo}}1 {ж}{{\cyrchar\cyrzh}}1 {з}{{\cyrchar\cyrz}}1
	{и}{{\cyrchar\cyri}}1 {й}{{\cyrchar\cyrishrt}}1 {к}{{\cyrchar\cyrk}}1
	{л}{{\cyrchar\cyrl}}1 {м}{{\cyrchar\cyrm}}1 {н}{{\cyrchar\cyrn}}1
	{о}{{\cyrchar\cyro}}1 {п}{{\cyrchar\cyrp}}1 {р}{{\cyrchar\cyrr}}1
	{с}{{\cyrchar\cyrs}}1 {т}{{\cyrchar\cyrt}}1 {у}{{\cyrchar\cyru}}1
	{ф}{{\cyrchar\cyrf}}1 {х}{{\cyrchar\cyrh}}1 {ц}{{\cyrchar\cyrc}}1
	{ч}{{\cyrchar\cyrch}}1 {ш}{{\cyrchar\cyrsh}}1 {щ}{{\cyrchar\cyrshch}}1
	{ъ}{{\cyrchar\cyrhrdsn}}1 {ы}{{\cyrchar\cyry}}1 {ь}{{\cyrchar\cyrsftsn}}1
	{э}{{\cyrchar\cyre}}1 {ю}{{\cyrchar\cyryu}}1 {я}{{\cyrchar\cyrya}}1
}
\usepackage{tocloft}
% Настройка отступов и переносов в оглавлении
\setlength{\cftsecindent}{1.5em} % Отступ для разделов
\setlength{\cftsubsecindent}{3.0em} % Отступ для подразделов
\setlength{\cftsecnumwidth}{2.5em} % Ширина номера раздела
\setlength{\cftsubsecnumwidth}{3.5em} % Ширина номера подраздела
% Разрешить перенос строк в оглавлении
\renewcommand{\cftsecfont}{\normalfont} % Обычный шрифт для заголовков
\renewcommand{\cftsubsecfont}{\normalfont}
\addto\extrasrussian{%
	\def\equationautorefname{формула}%
	\def\figureautorefname{рисунок}%
	\def\listingautorefname{листинг}%
	\def\tableautorefname{таблица}%
}

\addto\captionsrussian{
	\renewcommand\contentsname{\centerline{\bfseries\large{\MakeUppercase{Содержание}}}}
	\renewcommand{\bibsection}{\sectioncentered*{Список использованных источников}}
	\renewcommand{\lstlistingname}{Листинг}
}

% Список источников добавляется окружением thebibliography.
\usepackage{extsizes}


\sloppy

\usepackage[expansion=false]{microtype}
\newlength{\fivecharsapprox}
\setlength{\fivecharsapprox}{6ex}


\usepackage{indentfirst}
\setlength{\parindent}{\fivecharsapprox} 

\usepackage{lscape}
\usepackage[left=3cm,top=2.0cm,right=1.5cm,bottom=2.7cm]{geometry}

\frenchspacing

\usepackage{perpage}
\MakePerPage{footnote}

\makeatletter 
\def\@makefnmark{\hbox{\@textsuperscript{\normalfont\@thefnmark)}}}
\makeatother

\usepackage[bottom]{footmisc}

\makeatletter
\renewcommand{\thesection}{\arabic{section}}
\makeatother

\setcounter{secnumdepth}{3}

% Зачем: Настраивает отступ между таблицей с содержанимем и словом СОДЕРЖАНИЕ
\usepackage{tocloft}
\setlength{\cftbeforetoctitleskip}{-1em}
\setlength{\cftaftertoctitleskip}{1em}


% Зачем: Определяет отступы слева для записей в таблице содержания.
\makeatletter
\renewcommand{\l@section}{\@dottedtocline{1}{0.5em}{1.2em}}
\renewcommand{\l@subsection}{\@dottedtocline{2}{1.7em}{2.5em}}
\renewcommand{\l@subsubsection}{\@dottedtocline{3}{3.7em}{3.0em}}
\makeatother



% Зачем: Работа с колонтитулами
\usepackage{fancyhdr} 
\pagestyle{fancy}


\fancyhf{} 
\fancyfoot[R]{\thepage}
\renewcommand{\footrulewidth}{0pt} 
\renewcommand{\headrulewidth}{0pt}
\fancypagestyle{plain}{ 
	\fancyhf{}
	\rfoot{\thepage}}


% Зачем: Задает стиль заголовков раздела жирным шрифтом, прописными буквами, без точки в конце
\makeatletter
\renewcommand\section{%
	\clearpage\@startsection {section}{1}%
	{\fivecharsapprox}%
	{-1em \@plus -1ex \@minus -.2ex}%
	{1em \@plus .2ex}%
	{\raggedright\hyphenpenalty=10000\normalfont\large\bfseries\MakeUppercase}}
\makeatother


% Зачем: Задает стиль заголовков подразделов

\makeatletter
\renewcommand\subsection{
	\@startsection{subsection}{2}
	{\fivecharsapprox}
	{-1em \@plus -1ex \@minus -.2ex}
	{1em \@plus .2ex}
	{\raggedright\hyphenpenalty=10000\normalfont\normalsize\bfseries}}
\makeatother


% Зачем: Задает стиль заголовков пунктов
\makeatletter
\renewcommand\subsubsection{
	\@startsection{subsubsection}{3}%
	{\fivecharsapprox}%
	{-1em \@plus -1ex \@minus -.2ex}%
	{1em \@plus .2ex}
	{\raggedright\hyphenpenalty=10000\normalfont\normalsize\bfseries}}
\makeatother

% Зачем: для оформления введения и заключения, они должны быть выровнены по центру.
\makeatletter
\newcommand\sectioncentered{%
	\clearpage\@startsection {section}{1}%
	{\z@}%
	{-1em \@plus -1ex \@minus -.2ex}%
	{1em \@plus .2ex}%
	{\centering\hyphenpenalty=10000\normalfont\large\bfseries\MakeUppercase}%
}
\makeatother



% Зачем: Задает стиль библиографии
%\bibliographystyle{belarus-specific-utf8gost780u}
%\usepackage[final]{graphicx}
%\DeclareGraphicsExtensions{.png,.jpg}

% Зачем: Добавление подписей к рисункам
\usepackage{caption}
\usepackage{subcaption}

% Зачем: поворот ячеек таблиц на 90 градусов
\usepackage{rotating}
\DeclareRobustCommand{\povernut}[1]{\begin{sideways}{#1}\end{sideways}}

\DeclareRobustCommand{\x}[1]{\text{#1}}

% Зачем: Задание подписей, разделителя и нумерации частей рисунков
% сказали что надо без курсива
% \DeclareCaptionLabelFormat{stbfigure}{Рисунок \emph{#2}}
% \DeclareCaptionLabelFormat{stbtable}{Таблица \emph{#2}}
% \DeclareCaptionLabelFormat{stblisting}{Листинг \emph{#2}}
\DeclareCaptionLabelFormat{stbfigure}{Рисунок #2}
\DeclareCaptionLabelFormat{stbtable}{Таблица #2}
\DeclareCaptionLabelFormat{stblisting}{Листинг #2}
\DeclareCaptionLabelSeparator{stb}{~--~}
\captionsetup{labelsep=stb}
\captionsetup[figure]{labelformat=stbfigure, justification=centering, font={small}}
\captionsetup[lstlisting]{labelformat=stblisting,justification=centering, font={small}}
% сказали что надо слева
% \captionsetup[table]{labelformat=stbtable, justification=centering, font={small}}
\captionsetup[table]{labelformat=stbtable, justification=raggedright,singlelinecheck=false, font={small}}
\renewcommand{\thesubfigure}{\asbuk{subfigure}}

% Зачем: Окружения для оформления формул
\usepackage{calc}
\newlength{\lengthWordWhere}
\settowidth{\lengthWordWhere}{где}
\newenvironment{explanation}
{
	\begin{itemize}[leftmargin=0cm, itemindent=\lengthWordWhere + \labelsep , labelsep=\labelsep]
		
		\renewcommand\labelitemi{}
	}
	{
		\\[\parsep]
	\end{itemize}
}


\usepackage{tabularx}

\newenvironment{explanationx}
{
	\noindent 
	\tabularx{\textwidth}{@{}ll@{ --- } X }
}
{ 
	\\[\parsep]
	\endtabularx
}

\usepackage{amsmath}

\usepackage{amsfonts}
\usepackage{amssymb}
\usepackage{amsthm}
\usepackage{calc}
\usepackage{fp}
\usepackage{enumitem}

\makeatletter
\AddEnumerateCounter{\asbuk}{\@asbuk}{щ)}
\makeatother


\setlist{nolistsep}
\renewcommand{\labelenumi}{\asbuk{enumi})}
\renewcommand{\labelenumii}{\arabic{enumii})}


\setlist[itemize,0]{itemindent=\parindent + 2.2ex,leftmargin=0ex,label=--}
\setlist[enumerate,1]{itemindent=\parindent + 2.7ex,leftmargin=0ex}
\setlist[enumerate,2]{itemindent=\parindent + \parindent - 2.7ex}

% Зачем: Включение номера раздела в номер формулы. Нумерация формул внутри раздела.
\AtBeginDocument{\numberwithin{equation}{section}}

% Зачем: Включение номера раздела в номер таблицы. Нумерация таблиц внутри раздела.
\AtBeginDocument{\numberwithin{table}{section}}

% Зачем: Включение номера раздела в номер рисунка. Нумерация рисунков внутри раздела.
\AtBeginDocument{\numberwithin{figure}{section}}

% Зачем: Включение номера раздела в номер листинга. Нумерация листингов внутри раздела.
\AtBeginDocument{\numberwithin{lstlisting}{section}}


\usepackage{makecell}
\usepackage{multirow}
\usepackage{array}


\usepackage{textcomp}

\usepackage{siunitx}
\sisetup{
	binary-units = true,
	output-decimal-marker = {,},
	per-mode = symbol,
	range-phrase = --,
}
\DeclareSIUnit{\sample}{S}


%\newcommand{\ignore}[2e]%{\hspace{0in}#2}


\usepackage{verbatim}
\usepackage{xcolor}

%\AtBeginDocument{\numberwithin{lstlisting}{section}}

\usepackage[normalem]{ulem}

\renewcommand{\UrlFont}{\small\rmfamily\tt}

% Магия для подсчета разнообразных объектов в документе
\usepackage{lastpage}
\usepackage{totcount}
\regtotcounter{section}

\usepackage{etoolbox}

\newcounter{totfigures}
\newcounter{tottables}
\newcounter{totreferences}
\newcounter{totequation}

\providecommand\totfig{} 
\providecommand\tottab{}
\providecommand\totref{}
\providecommand\toteq{}

\makeatletter
\AtEndDocument{%
	\addtocounter{totfigures}{\value{figure}}%
	\addtocounter{tottables}{\value{table}}%
	\addtocounter{totequation}{\value{equation}}
	\immediate\write\@mainaux{%
		\string\gdef\string\totfig{\number\value{totfigures}}%
		\string\gdef\string\tottab{\number\value{tottables}}%
		\string\gdef\string\totref{\number\value{totreferences}}%
		\string\gdef\string\toteq{\number\value{totequation}}%
	}%
}
\makeatother

\pretocmd{\section}{\addtocounter{totfigures}{\value{figure}}\setcounter{figure}{0}}{}{}
\pretocmd{\section}{\addtocounter{tottables}{\value{table}}\setcounter{table}{0}}{}{}
\pretocmd{\section}{\addtocounter{totequation}{\value{equation}}\setcounter{equation}{0}}{}{}
\pretocmd{\bibitem}{\addtocounter{totreferences}{1}}{}{}



% Для оформления таблиц не влязящих на 1 страницу
\usepackage{longtable}

\usepackage{gensymb}

% Зачем: преобразовывать текст в верхний регистр командой MakeTextUppercase
\usepackage{textcase}

%  Переносы в словах с тире \hyph.

\def\hyph{-\penalty0\hskip0pt\relax}

% Добавляем левый отступ для библиографии

\makeatletter
\renewenvironment{thebibliography}[1]
{\sectioncentered*{Список использованных источников}
	\@mkboth{\MakeUppercase\refname}{\MakeUppercase\refname}%
	\list{\@biblabel{\@arabic\c@enumiv}}%
	{\settowidth\labelwidth{\@biblabel{#1}}%
		\setlength{\itemindent}{\dimexpr\labelwidth+\labelsep+1em}
		\leftmargin\z@
		\@openbib@code
		\usecounter{enumiv}%
		\let\p@enumiv\@empty
		\renewcommand\theenumiv{\@arabic\c@enumiv}}%
	\sloppy
	\clubpenalty4000
	\@clubpenalty \clubpenalty
	\widowpenalty4000%
	\sfcode`\.\@m}
{\def\@noitemerr
	{\@latex@warning{Empty `thebibliography' environment}}%
	\endlist}
\makeatother

\newcommand{\intro}[3]{
	\stepcounter{section}
	\sectioncentered*{ПРИЛОЖЕНИЕ \MakeUppercase{#1}}
	\begin{center} 
		\bf{(#2)}\\
		\bf{#3}
	\end{center}
	\markboth{\MakeUppercase{#1}}{}
	\addcontentsline{toc}{section}{Приложение \MakeUppercase{#1} (#2) #3}
}
\usepackage{enumitem}
\usepackage{url} 

\usepackage{float}
\usepackage{amsmath}

\newcommand{\bnfterm}[1]{\langle\text{#1}\rangle}
\newcommand{\bnfdef}{::=}
\newcommand{\bnfor}{\mid}

% Для длинных определений
\newenvironment{longbnf}
{\begin{tabular}{@{}l@{}}}
	{\end{tabular}}
\usepackage{tabularx} % в преамбулу

\setcounter{tocdepth}{3}  % Включает в оглавление section и subsection
% В preamble.tex найдите место, где подключаете пакеты
